%%%%%%%%%%%%%%%%%%%%%%%%%%%%%%%%%%%%%%%%%
% Awesome Cover Letter
% XeLaTeX Template
% Version 1.1 (9/1/2016)
%
% This template has been downloaded from:
% http://www.LaTeXTemplates.com
%
% Original authors:
% Claud D. Park (posquit0.bj@gmail.com)
% Lars Richter (mail@ayeks.de)
% With modifications by:
% Vel (vel@latextemplates.com)
%
% License:
% CC BY-NC-SA 3.0 (http://creativecommons.org/licenses/by-nc-sa/3.0/)
%
% Important note:
% This template must be compiled with XeLaTeX, the below lines will ensure this
%!TEX TS-program = xelatex
%!TEX encoding = UTF-8 Unicode
%
%%%%%%%%%%%%%%%%%%%%%%%%%%%%%%%%%%%%%%%%%

%----------------------------------------------------------------------------------------
%	PACKAGES AND OTHER DOCUMENT CONFIGURATIONS
%----------------------------------------------------------------------------------------

\documentclass[10pt, letter]{ag-cv} % A4 paper size by default, use 'letterpaper' for US letter

\geometry{left=2cm, top=1.3cm, right=2cm, bottom=1.9cm, footskip=.5cm} % Configure page margins with geometry
 
\fontdir[fonts/] % Specify the location of the included fonts
\usepackage{ragged2e}

% Color for highlights
\definecolor{awesome}{HTML}{1c85c7} % Default colors include: awesome-emerald, awesome-skyblue, awesome-red, awesome-pink, awesome-orange, awesome-nephritis, awesome-concrete, awesome-darknight
%\definecolor{awesome}{HTML}{CA63A8} % Uncomment if you would like to specify your own color

% Colors for text - uncomment and modify
%\definecolor{darktext}{HTML}{414141}
%\definecolor{text}{HTML}{414141}
%\definecolor{graytext}{HTML}{414141}
%\definecolor{lighttext}{HTML}{414141}

\renewcommand{\acvHeaderSocialSep}{\quad\textbar\quad} % If you would like to change the social information separator from a pipe (|) to something else

%----------------------------------------------------------------------------------------
%	PERSONAL INFORMATION
%	Comment any of the lines below if they are not required
%----------------------------------------------------------------------------------------



\name{}{Alvaro Gomariz}
\address{B\"andlistrasse 45, 
CH-8064 Zürich}
% \mobile{(+82) 10-9030-1843}

\email{alvaro.gomariz@roche.com}
\orcid{0000-0002-6172-5190}
% \github{https://github.com/alvgom}
\linkedin{alvarogomariz}
\gscholar{https://scholar.google.com/citations?hl=en&user=v4JtSZcAAAAJ}
\mobile{ +41 (0)76 773 14 72}


% \position{Software Engineer{\enskip\cdotp\enskip}Security Expert} % Your expertise/fields
% \quote{``Make the change that you want to see in the world."} % A quote or statement

%----------------------------------------------------------------------------------------
%	RECIPIENT/POSITION/LETTER INFORMATION
%	All of the below lines must be filled out
%----------------------------------------------------------------------------------------

% \recipient{}{Roche} % The company being applied to

\letterdate{Zurich, 28th February 2022} % The date on the letter, default is the date of compilation

\lettertitle{Application for Machine Learning Scientist - gRED} % The title of the letter

\letteropening{Dear Dr. Loukas,} % How the letter is opened

\letterclosing{Yours sincerely,} % How the letter is closed

% \letterenclosure[Attached]{Curriculum Vitae} % Any enclosures with the letter

% \makecvfooter{\today}{Alvaro Gomariz~~~·~~~Cover Letter}{} % Specify the letter footer with 3 arguments: (<left>, <center>, <right>), leave any of these blank if they are not needed
  
%----------------------------------------------------------------------------------------

\begin{document}

\makecvheader % Print the header

\makelettertitlenorec % Print the title

%----------------------------------------------------------------------------------------
%	LETTER CONTENT
%----------------------------------------------------------------------------------------

\begin{cvletter}

%------------------------------------------------

% \lettersection{About Me}




% \lettersection{Why Me?}
\begin{flushleft}
Advancing a data-based understanding of healthcare has been the passion inspiring my career. Working at Roche as a postdoc in the past year, I have been able to follow this vision, and I am eager to continue contributing to this goal in the long-term. It is for this reason that I could not be more excited to find this position, as it perfectly aligns with my interests and skills. 

The beginning of my PhD at the Computer Vision Lab at ETH Zurich, and in collaboration with the University Hospital Zurich, coincided with the surge in deep learning research. This was key for my project, which was focused on the analysis of complex multidimensional microscopy datasets. I hence acquired a solid understanding of the current machine learning (ML) literature and made my own contributions; the results of my research have been published in renowned scientific journals, including Nature Communications, Nature Machine Intelligence, and Science Advances. My work has also been recognized with the best PhD Award 2021 in the field of biomedical imaging by the EXCITE Zurich centre. 

My interest in artificial intelligence (AI) research always went beyond the design of its tools, as I was fascinated by the application possibilities to numerous biomedical prediction tasks; over the course of my PhD I have initiated several collaborations with biologists, clinicians, and healthcare companies to apply my research to different image modalities and underlying questions. I also applied my ML expertise as an independent consultant for designing and delivering a product for image segmentation and spatial statistics to clients in the biotechnology sector, where I gained experience in supervising a team of software engineers.  
 
I have always been passionate about the quest to advance healthcare, which is why I chose to begin my current postdoc project in pRED informatics at Roche. In this position, I have learned to work following Agile development principles and design new deep learning solutions in order to advance the analysis capabilities and reduce the manual annotation efforts in PHC ophthalmology. 

I am thus very excited about the open position as Machine Learning Scientist in gRED because it would allow me to employ my knowledge in applied AI in biology and lead to impactful research and computational solutions. While several ML foundations are similar, the focus of my career so far has been on computer vision. Hence, I am very excited to expand my ML skills in the fields of graph representation learning, geometric deep learning, and reinforcement learning.  

I believe that my experience and accomplishments make me an excellent fit for this position, as they have allowed me to develop the technical expertise working with other complex biomedical datasets, an understanding of the underlying needs of the analyses, as well as the skills for effective communication and collaboration across disciplines to bridge the gap between computer science and biology / medicine. 

Thank you for taking the time to review my application. I am highly motivated to apply and expand my ML expertise to graph and sequence-based biological data, and I trust that I can bring new and creative ideas for impactful findings.

\end{flushleft}

\end{cvletter}

%----------------------------------------------------------------------------------------

\makeletterclosing % Print the signature and enclosures

\end{document}